\scanpage{12}
Here is a candidate for a functor, which may \unsure{tumn|turn} out to be representable, in which case this would be $P$. Let $\underline{A}$ be the ring-scheme over $k$ defined by $A$ (see Greenberg's papers, or Serre's in Bulletin Soc. Math.), so that
\[
A\handwritten{\simeq}\underline{A}(k)\ .
\tag*{(23)}
\]
For every algebra $B$ over $k$, corsider the ring $\underline{A}(B)$, which is an algebra over $\underline{A}(k)=A$. We may thus consider
\[
X_B\overset{\mathrm{dfn}}{=}X\handwritten{\otimes}_A\underline{A}(B)\ ,
\]
and consider the functor
\[
B\handwritten{\longmapsto}\operatorname{Pic}(X_B)\ .
\tag*{(\unsure{34|24})}
\]
This of course may not be representable, even if $A=k$. However, it may turn out that the fppf sheaf associated to this functor is representable, and this then would give the looked for cand\deleted{\unsure{r}}idate for $P$. In any case, Artin's work gives us a very handy set of necessary and sufficient conditions for a group functor to be representable by a scheme locally of finite type, and in principle it should be possible to decide wether or not the \deleted{functor} sheaf just considered satisfies to these conditions or not. If so, one may expect that the developments on local Picard schemes given for the case when there is a field of representatives will carry over to the case when the residue field is perfect with car.$\handwritten{>}0$.
% archive sheet number: 11
